\section{Robust Anytime-Valid Detection under Static Reference Contamination}

\subsection{Candidate-wise evidence}

For each \(C\in\mathcal C_m\), apply the same pre-specified clean predictive-rank betting functional to the reconstructed sequence
\(R_1^{(C)},R_2^{(C)},\ldots,\)
producing candidate wealth \(E_t^{(C)}\).

For \(C\neq C^\star\), the reconstructed sequence need not obey the clean predictive law, so \(E_t^{(C)}\) is not claimed to be a martingale. For the true candidate, Lemma~1 and \eqref{eq:canonical-clean-rank} give
\(R_t^{(C^\star)}=R_t^\star.\)
Because the betting rule is a deterministic predictable functional of the candidate rank history, the true-candidate bet \(q_t^{(C^\star)}\) is \(\mathcal G_{t-1}^\star\)-measurable. Therefore \(E_t^{(C^\star)}\) is exactly the clean PRM with respect to the canonical clean-rank filtration \((\mathcal G_t^\star)\).

The robust evidence lower envelope is defined by Equation~\eqref{eq:lower-envelope}:
\begin{equation}
  \underline E_t
  =
  \min_{C\in\mathcal C_m}E_t^{(C)}.
  \label{eq:lower-envelope}
\end{equation}
Since \(C^\star\in\mathcal C_m\), the pathwise domination relation is given by Equation~\eqref{eq:pathwise-domination}:
\begin{equation}
  \underline E_t
  \le
  E_t^{(C^\star)}
  \qquad\text{for every }t
  \label{eq:pathwise-domination}
\end{equation}
The domination holds pathwise.

\begin{lemma}[Random sorted positions do not create post-selection]
Let \(\Phi_t\) denote the deterministic wealth functional induced by the pre-specified betting strategy on a clean rank history of length \(t\). Equation~\eqref{eq:true-candidate-functional} expresses the true-candidate wealth as the same functional of the canonical clean ranks:
\begin{equation}
E_t^{(C^\star)}
  =
  \Phi_t(R_1^\star,\ldots,R_t^\star)
  \label{eq:true-candidate-functional}
\end{equation}
This identity holds pathwise, even though the sorted position set \(C^\star\) is random.
\end{lemma}

\begin{proof}
For every realized reference bank, deleting the realized contaminated positions \(C^\star\) gives the canonical clean ranks by Lemma~1. The candidate algorithm then applies the same deterministic betting functional \(\Phi_t\) to those ranks. Thus the process indexed by the random representation \(C^\star\) is identically the single canonical clean-rank wealth process; it is not chosen from unrelated martingales after observing online evidence.
\end{proof}

\begin{theorem}[Anytime-valid robust crossing under exactly \(m\) static contaminants]
Assume the model of Section~2 with exactly \(m\) static/exogenous contaminated references. Let \(E_t^{(C)}\) be the candidate wealth obtained by applying the same pre-specified valid clean-reference betting functional to each reconstructed candidate sequence. Then, for every \(\alpha\in(0,1)\), the anytime-valid crossing bound in Equation~\eqref{eq:main-theorem} holds:
\begin{equation}
  \Pr_{H_0}\!\left(
    \sup_{t\ge0}\underline E_t\ge\frac1\alpha
  \right)
  \le \alpha.
  \label{eq:main-theorem}
\end{equation}
\end{theorem}

\begin{proof}
The preceding lemma identifies \(E_t^{(C^\star)}\) with the canonical clean PRM
\(\Phi_t(R_1^\star,\ldots,R_t^\star).\)
Hence \(E_t^{(C^\star)}\) is a nonnegative martingale with respect to \((\mathcal G_t^\star)\). By Equation~\eqref{eq:pathwise-domination}, the crossing-event inclusion in Equation~\eqref{eq:crossing-event-inclusion} follows:
\begin{equation}
\begin{aligned}
\left\{\sup_t\underline E_t\ge1/\alpha\right\}
&\subseteq\\
\left\{\sup_tE_t^{(C^\star)}\ge1/\alpha\right\}.
\end{aligned}
\label{eq:crossing-event-inclusion}
\end{equation}
Ville's inequality for the canonical true-candidate martingale yields \eqref{eq:main-theorem}.
\end{proof}

The theorem establishes anytime-valid \emph{threshold-crossing control} for \(\underline E_t\). It does not claim that \((\underline E_t)\) itself is a martingale, supermartingale, or e-process.

\subsection{Unknown count bounded by M}

The exact value of the contamination count is not needed for validity if a deterministic upper bound is available.

For each
\(r\in\{0,\ldots,M\},\)
Equation~\eqref{eq:bounded-contamination-candidates} defines the candidate sets and number of clean references:
\begin{equation}
\begin{aligned}
\mathcal C_r
&=\{C\subseteq\{1,\ldots,K\}:|C|=r\},\\
n_r&=K-r.
\end{aligned}
\label{eq:bounded-contamination-candidates}
\end{equation}
For \(C\in\mathcal C_r\), the rank reconstruction is given by Equation~\eqref{eq:bounded-contamination-rank}:
\begin{equation}
R_t^{(r,C)}
  =
  Q_t-\#\{c\in C:c\le Q_t\},
  \label{eq:bounded-contamination-rank}
\end{equation}
We then compute a valid clean-reference predictive-rank wealth \(E_t^{(r,C)}\) using the null law corresponding to \(n_r\) clean references.

The upper-bound envelope is defined by Equation~\eqref{eq:upper-bound-envelope}:
\begin{equation}
  \underline E_t^{(\le M)}
  =
  \min_{0\le r\le M}
  \min_{C\in\mathcal C_r}
  E_t^{(r,C)}.
  \label{eq:upper-bound-envelope}
\end{equation}

\begin{corollary}[Anytime validity with only an upper bound]
Suppose the true number of static/exogenous contaminated references is \(m^\star\le M\). Then, for every \(\alpha\in(0,1)\), the anytime-valid crossing bound in Equation~\eqref{eq:upper-bound-validity} holds:
\begin{equation}
  \Pr_{H_0}\!\left(
    \sup_{t\ge0}
    \underline E_t^{(\le M)}
    \ge\frac1\alpha
  \right)
  \le\alpha.
  \label{eq:upper-bound-validity}
\end{equation}
\end{corollary}

\begin{proof}
The candidate family in \eqref{eq:upper-bound-envelope} contains the true pair \((m^\star,C^\star)\). For that pair, \(n_{m^\star}=K-m^\star\), the reconstructed ranks equal the canonical clean ranks, and \(E_t^{(m^\star,C^\star)}\) is a valid clean PRM. Therefore
\(\underline E_t^{(\le M)} \le E_t^{(m^\star,C^\star)}\)
pathwise for all \(t\). Ville's inequality gives \eqref{eq:upper-bound-validity}.
\end{proof}

This corollary is a theoretical extension of the validity result. The frozen experiments still use the exact count \(m=2\); no empirical claim for the unknown-count procedure is made here.

\subsection{Persistent partition structure}

The observable rank \(Q_t\) takes values in \(\{0,\ldots,K\}\), defining \(K+1\) adjacent observed-rank bins. A sorted reference position \(c\in\{1,\ldots,K\}\) is precisely the boundary between observed-rank bins \(c-1\) and \(c\).

\begin{proposition}[Bijection between candidates and contiguous partitions]
For fixed \(m\), candidate sets
\(C\subseteq\{1,\ldots,K\}, \qquad |C|=m,\)
are in one-to-one correspondence with contiguous partitions of the \(K+1\) ordered observed-rank bins into
\(K-m+1=n+1\)
nonempty blocks.
\end{proposition}

\begin{proof}
Deleting a candidate reference at position \(c\) removes exactly the boundary between bins \(c-1\) and \(c\), merging the two connected components adjacent to that boundary. Deleting \(m\) distinct boundaries from the path of \(K+1\) ordered bins leaves exactly \(K+1-m=n+1\) connected components, each a nonempty contiguous block.

Conversely, any partition of the \(K+1\) ordered bins into \(n+1\) nonempty contiguous blocks has exactly \(m\) missing internal boundaries. Taking those missing boundary indices as \(C\) recovers a unique candidate set. The two maps are inverses.
\end{proof}

The same partition is reused for all \(t\), which is the combinatorial expression of persistent contamination identities.

\subsection{Fixed-categorical closed form}

Consider a fixed categorical betting distribution
\(a=(a_0,\ldots,a_n), \qquad a_j>0,\qquad \sum_{j=0}^n a_j=1.\)
The full-reference rank histogram is defined by Equation~\eqref{eq:observed-rank-histogram}:
\begin{equation}
\begin{aligned}
H_q(t)&=\sum_{s=1}^t\mathbf 1\{Q_s=q\},\\
&\qquad q=0,\ldots,K.
\end{aligned}
\label{eq:observed-rank-histogram}
\end{equation}
Let \(N_j^{(C)}(t)\) be the histogram after the candidate merges. Equation~\eqref{eq:fixed-categorical-closed} gives the candidate log wealth:
\begin{align}
  \log E_t^{(C)}
  &=
  \log\Gamma(n+t+1) \nonumber\\
  &\quad-\log\Gamma(n+1) \nonumber\\
  &\quad+
  \sum_{j=0}^n
  \left[
    N_j^{(C)}(t)\log a_j
    \right. \nonumber\\
  &\qquad\left.
    -
    \log\Gamma(N_j^{(C)}(t)+1)
  \right].
  \label{eq:fixed-categorical-closed}
\end{align}
Thus the candidate objective depends on the online history through the observed-rank histogram and the contiguous partition induced by \(C\).

\begin{theorem}[Exact dynamic program for a fixed categorical bettor]
After \(O(K)\) prefix-sum preprocessing of the observed-rank histogram, the minimum of \eqref{eq:fixed-categorical-closed} over all \(C\in\mathcal C_m\) can be computed in
\(O(Km^2)\)
time and
\(O(Km)\)
memory.
\end{theorem}

\begin{proof}
Let the ordered observed-rank bins be indexed \(0,\ldots,K\). Equation~\eqref{eq:segment-histogram-count} defines the total histogram count in a segment \([b,b+\ell-1]\), evaluable in \(O(1)\) time after prefix-sum preprocessing:
\begin{equation}
N(b,\ell)
  =
  \sum_{q=b}^{b+\ell-1}H_q(t).
  \label{eq:segment-histogram-count}
\end{equation}
If this segment is the next contiguous block and is assigned to clean-rank index \(j\), its additive contribution to Equation~\eqref{eq:fixed-categorical-closed}, excluding constants independent of the partition, is given by Equation~\eqref{eq:segment-additive-cost}:
\begin{equation}
\begin{aligned}
c(j,b,\ell)
&=N(b,\ell)\log a_j\\
&\quad-\log\Gamma(N(b,\ell)+1).
\end{aligned}
\label{eq:segment-additive-cost}
\end{equation}

Define \(DP[b,d]\) as the minimum additive cost among all valid contiguous partitions of the first \(b\) observed-rank bins that have deleted exactly \(d\) internal boundaries. Such a partial partition contains exactly \(b-d\) completed blocks, so the next block receives clean-rank index
\(j=b-d.\)
Choosing its length \(\ell\ge1\) deletes exactly \(\ell-1\) additional boundaries. Set \(b'=b+\ell\) and \(d'=d+\ell-1\); the dynamic-programming update in Equation~\eqref{eq:dp-recurrence} is algebraically the same recurrence:
\begin{equation}
\begin{aligned}
\mathrm{DP}[b',d'] &\gets
  \min\!\bigl\{\mathrm{DP}[b',d'],\\
&\qquad \mathrm{DP}[b,d]+c(b-d,b,\ell)\bigr\}.
\end{aligned}
\label{eq:dp-recurrence}
\end{equation}
Only transitions satisfying
\(d+\ell-1\le m\)
and \(b+\ell\le K+1\) are allowed.

By the preceding bijection, every candidate set corresponds to exactly one sequence of contiguous blocks and hence to exactly one path from \(DP[0,0]\) to \(DP[K+1,m]\). Conversely, every feasible path in the recurrence specifies a unique contiguous partition and therefore a unique candidate. Induction on \(b\) shows that each state stores the minimum cost over all and only feasible partial partitions represented by that state. Thus \(DP[K+1,m]\) is the exact candidate minimum.

There are \(O(Km)\) states. From each state, at most \(m+1\) segment lengths can be feasible because \(\ell-1\le m-d\). Segment costs are \(O(1)\) after prefix-sum preprocessing. Hence the total running time is \(O(Km^2)\). Storing the table and optional predecessors requires \(O(Km)\) memory.
\end{proof}

The frozen empirical method uses a convex portfolio of wealth processes. The logarithm of that mixture contains a log-sum-exp across components, so the additive structure above does not directly establish the same polynomial-time result for the complete frozen portfolio. We therefore restrict the computational theorem to the fixed-categorical case.

\subsection{Finite-reference betting asymptotics}

The finite-reference information ceiling can be made explicit for the two betting families used in the frozen portfolio. The following proposition is an asymptotic characterization of these components, not a claim of a new general Bayes-factor principle.

Equation~\eqref{eq:empirical-rank-frequencies} defines the category count and empirical rank frequencies:
\begin{equation}
\begin{aligned}
k&=n+1,\\
\widehat\pi_{j,t}&=\frac{N_j(t)}{t}.
\end{aligned}
\label{eq:empirical-rank-frequencies}
\end{equation}
Assume the interior-frequency condition in Equation~\eqref{eq:interior-rank-frequencies} holds:
\begin{equation}
  \widehat\pi_{j,t}\longrightarrow\pi_j>0,
  \qquad j=0,\ldots,n.
  \label{eq:interior-rank-frequencies}
\end{equation}

\begin{proposition}[Finite-reference asymptotics]
Under \eqref{eq:interior-rank-frequencies}:

\begin{enumerate}[label=(\alph*),leftmargin=*]
\item For a fixed categorical bettor \(a=(a_0,\ldots,a_n)\) with all \(a_j>0\), the limiting per-time log wealth is stated in Equation~\eqref{eq:fixed-cat-asymptotic}:
\begin{equation}
  \frac1t\log E_t^{(a)}
  \longrightarrow
  -D_{\mathrm{KL}}(\pi\Vert a).
  \label{eq:fixed-cat-asymptotic}
\end{equation}

\item For a Dirichlet predictive bettor with parameters
\(\boldsymbol\alpha=(\alpha_0,\ldots,\alpha_n), \qquad \alpha_j>0, \qquad A=\sum_{j=0}^n\alpha_j,\)
whose predictive probabilities are given by Equation~\eqref{eq:dirichlet-betting-probabilities}:
\begin{equation}
q_{t,j}^{(\boldsymbol\alpha)}
  =
  \frac{N_{j,t-1}+\alpha_j}{t-1+A},
  \label{eq:dirichlet-betting-probabilities}
\end{equation}
The wealth has the exact form in Equation~\eqref{eq:dirichlet-wealth-exact}:
\begin{equation}
\begin{aligned}
E_t^{(\boldsymbol\alpha)}
&=\frac{\Gamma(A)\Gamma(k+t)}
        {\Gamma(A+t)\Gamma(k)}\\
&\quad\times\prod_{j=0}^n
  \frac{\Gamma(N_j(t)+\alpha_j)}
       {\Gamma(\alpha_j)\Gamma(N_j(t)+1)}.
\end{aligned}
\label{eq:dirichlet-wealth-exact}
\end{equation}
It converges to the finite positive limit in Equation~\eqref{eq:dirichlet-wealth-limit}:
\begin{equation}
  E_t^{(\boldsymbol\alpha)}
  \longrightarrow
  \frac{\Gamma(A)}
       {\Gamma(k)\prod_{j=0}^n\Gamma(\alpha_j)}
  \prod_{j=0}^n
  \pi_j^{\alpha_j-1}.
  \label{eq:dirichlet-wealth-limit}
\end{equation}
\end{enumerate}
\end{proposition}

\begin{proof}
For part (a), Equation~\eqref{eq:log-wealth-expansion} expands the fixed-categorical log wealth:
\begin{equation}
\begin{aligned}
  \frac1t\log E_t^{(a)}
  &=
  \frac{\log\Gamma(t+k)}{t}\\
  &\quad-
  \frac1t\sum_{j=0}^n\log\Gamma(N_j(t)+1)\\
  &\quad+
  \sum_{j=0}^n\widehat\pi_{j,t}\log a_j
  +o(1).
\end{aligned}
  \label{eq:log-wealth-expansion}
\end{equation}
Stirling's formula and \(\widehat\pi_{j,t}\to\pi_j>0\) imply the limit in Equation~\eqref{eq:stirling-entropy-limit}:
\begin{equation}
\begin{aligned}
&\frac{\log\Gamma(t+k)}{t}\\
&\quad-\frac1t\sum_{j=0}^n\log\Gamma(N_j(t)+1)\\
&\longrightarrow -\sum_{j=0}^n\pi_j\log\pi_j.
\end{aligned}
\label{eq:stirling-entropy-limit}
\end{equation}
Adding \(\sum_j\pi_j\log a_j\) yields
\(-\sum_j\pi_j\log\frac{\pi_j}{a_j} = -D_{\mathrm{KL}}(\pi\Vert a).\)

For part (b), multiplying the Dirichlet predictive factors telescopes into the ratio of two Dirichlet--multinomial sequence probabilities, yielding \eqref{eq:dirichlet-wealth-exact}. The gamma-ratio asymptotic
\(\frac{\Gamma(x+\alpha)}{\Gamma(x+\beta)} \sim x^{\alpha-\beta}\)
gives
\(\frac{\Gamma(k+t)}{\Gamma(A+t)} \sim t^{k-A}\)
and, for each \(j\), the count-dependent asymptotic in Equation~\eqref{eq:dirichlet-count-asymptotic}:
\begin{equation}
\begin{aligned}
\frac{\Gamma(N_j(t)+\alpha_j)}
     {\Gamma(N_j(t)+1)}
&\sim N_j(t)^{\alpha_j-1}\\
&\sim t^{\alpha_j-1}\pi_j^{\alpha_j-1}.
\end{aligned}
\label{eq:dirichlet-count-asymptotic}
\end{equation}
Since
\(\sum_{j=0}^n(\alpha_j-1)=A-k,\)
the powers of \(t\) cancel, leaving \eqref{eq:dirichlet-wealth-limit}.
\end{proof}

A direct consequence is that if the limiting interior rank-frequency vector \(\pi\) differs from every fixed categorical vector in a finite betting portfolio, then every fixed categorical component decays exponentially by \eqref{eq:fixed-cat-asymptotic}, while every Dirichlet predictive component converges to a finite constant by \eqref{eq:dirichlet-wealth-limit}. Adding a constant-one hedge therefore yields a finite limiting true-candidate portfolio wealth. Since
\(\underline E_t\le E_t^{(C^\star)},\)
the robust lower envelope cannot diverge to \(+\infty\) in such a case.

This provides a formal explanation for the finite-reference information ceiling: with \(n\) fixed, the frozen portfolio is not expected to deliver universal asymptotic wealth divergence across all persistent rank-frequency alternatives. The paper therefore focuses on finite-horizon anytime validity and finite-reference power.

\subsection{Scope boundary}

The theory still excludes value-adaptive replacement after observing an originally clean reference sample. In that stronger model the surviving clean observations can be selection-biased and the canonical clean-rank PRM need not apply.

The upper-bound corollary extends validity from known \(m\) to known \(M\), but the current experimental study remains frozen at exact \(m=2\). Likewise, the asymptotic proposition explains the frozen bettor family but does not imply an optimality statement for other betting strategies.
